\section{Glossar}
\begin{description}[style=nextline, itemsep=2pt]
  \item[ACID] \emph{Atomicity, Consistency, Isolation, Durability} -- die vier Eigenschaften zuverlässiger Transaktionen.
  \item[Aggregatfunktion] Funktion, die mehrere Zeilen zu einem einzelnen Wert zusammenfasst: \icode{COUNT}, \icode{SUM}, \icode{AVG}, \icode{MIN}, \icode{MAX}.
  \item[Anomalie] Fehler, der durch Redundanz in unnormalisierten Tabellen entsteht (Einfüge-, Änderungs- und Löschanomalie).
  \item[ANSI-3-Ebenen-Modell] Architektur, die eine Datenbank in eine interne, eine konzeptionelle und eine externe Ebene trennt (auch \emph{SPARC} genannt).
  \item[Attribut] Eigenschaft eines Objekts bzw. Spalte einer Relation.
  \item[CAP-Theorem] Konsistenz, Verfügbarkeit und Partitionstoleranz sind in einem verteilten System nie gleichzeitig vollständig erreichbar -- nur zwei davon.
  \item[Caching] Clients halten eine eigene, lokale Kopie der Daten (oder von Teilen davon) -- schneller Zugriff, aber Konsistenzfragen.
  \item[Chen-Notation] Standardnotation für ER-Diagramme: Rechteck = Objekttyp, Ellipse = Eigenschaft, Raute = Beziehungstyp.
  \item[Constraint] Regel auf einer Spalte oder Tabelle, z.\,B. \icode{NOT NULL} oder \icode{UNIQUE}.
  \item[CQRS] \emph{Command Query Responsibility Segregation} -- Architekturmuster, das Schreiben (Commands, Write-Modell) und Lesen (Queries, Read-Modell) strikt trennt; Synchronisation über Events.
  \item[Data Mart] Abteilungs- oder funktionsbezogene Teilmenge eines Data Warehouse.
  \item[Data Mining] Automatisches Entdecken von Modellen und Mustern in großen Datenbeständen.
  \item[Data Warehouse] Sehr große, von den operativen Systemen unabhängige Datenbank mit verdichteten Daten für Auswertungen.
  \item[Datenwürfel (Cube)] Mehrdimensionale Sicht auf Kennzahlen: Jede Zelle enthält eine Kennzahl am Schnittpunkt der Dimensionen.
  \item[DBMS] Datenbank-Management-System -- verwaltet die Datenbank und ist die Schnittstelle zu ihr (z.\,B. MySQL, PostgreSQL).
  \item[DCL] \emph{Data Control Language} -- steuert Zugriff und Rechte (\icode{GRANT}, \icode{REVOKE}).
  \item[DDL] \emph{Data Definition Language} -- definiert und verändert das Schema (\icode{CREATE}, \icode{ALTER}, \icode{DROP}).
  \item[Deadlock] Gegenseitiges Sperren zweier Transaktionen, die aufeinander warten.
  \item[Dimension] Blickwinkel, aus dem eine Kennzahl (Fakt) betrachtet wird -- z.\,B. Zeit, Produkt, Kunde oder Region; besitzt Hierarchien (Tag $\rightarrow$ Monat $\rightarrow$ Jahr).
  \item[DML] \emph{Data Manipulation Language} -- verändert Daten in bestehenden Tabellen (\icode{INSERT}, \icode{UPDATE}, \icode{DELETE}).
  \item[DQL] \emph{Data Query Language} -- fragt Daten ab (\icode{SELECT}).
  \item[DTO] \emph{Data Transfer Object} -- Objekt, das Daten zwischen Anwendung und API transportiert; Teil der externen Ebene.
  \item[Entität] Konkretes Objekt der realen Welt, das in der Datenbank modelliert wird.
  \item[ERM] Entity-Relationship-Modell -- konzeptioneller Entwurf einer Datenbank mit Entitäten, Attributen und Beziehungen.
  \item[Event] Automatisch geplante SQL-Anweisung, die zu einer bestimmten Zeit oder regelmäßig ausgeführt wird.
  \item[Foreign Key] Fremdschlüssel -- verweist auf den Primary Key einer anderen Tabelle und verknüpft so die Daten.
  \item[Funktionale Abhängigkeit] $A \rightarrow B$: Zu jedem Wert von $A$ gehört genau ein Wert von $B$.
  \item[Isolationsebene] Legt fest, wie stark gleichzeitige Transaktionen getrennt werden: \icode{READ UNCOMMITTED} bis \icode{SERIALIZABLE}.
  \item[Join] Verknüpft Daten aus mehreren Tabellen über Primär- und Fremdschlüssel (\icode{INNER}, \icode{LEFT}, \icode{RIGHT}, \icode{CROSS}).
  \item[Kardinalität] Gibt an, wie viele Objekte eines Typs an einer Beziehung teilnehmen können (1:1, 1:n, n:m).
  \item[Locking] Temporäres Sperren von Datenbankobjekten -- \emph{Shared Lock} zum Lesen, \emph{Exclusive Lock} zum Schreiben.
  \item[Managementunterstützungssystem] System zur Extraktion von Wissen aus großen Datenmengen für Berichte über vergangene Daten -- Basis: Data Warehouse, OLAP und Data Mining.
  \item[Map/Reduce] Programmiermodell zur Parallelisierung nach dem Prinzip \emph{Divide and Conquer}.
  \item[Multiplizität] Zahlenverhältnis zwischen Objekten im UML-Klassendiagramm, z.\,B. \icode{1}, \icode{0..5}, \icode{0..*}.
  \item[Normalisierung] Schrittweises Aufteilen von Relationen (1NF--3NF), um Redundanz und Anomalien zu vermeiden.
  \item[NoSQL] \emph{Not Only SQL} -- Alternative zu relationalen DBMS (Key-Value, Document Store, Column Store, Graphdatenbank).
  \item[OLAP] \emph{Online Analytical Processing} -- flexible, mehrdimensionale Datenanalyse nach dem Top-down-Ansatz.
  \item[OLTP] \emph{Online Transaction Processing} -- Verarbeitung von Transaktionen in der Anwendungsdatenbank; schreiboptimiert und normalisiert.
  \item[ORM] \emph{Object-Relational Mapper} -- bildet Objekte einer Anwendung auf relationale Tabellen ab (z.\,B. Entity Framework Core).
  \item[Primary Key] Primärschlüssel -- identifiziert einen Datensatz eindeutig, ist nie \icode{NULL} und kann aus mehreren Spalten bestehen.
  \item[Prozedur] Benannter, in der Datenbank gespeicherter Anweisungsblock; Aufruf per \icode{CALL}.
  \item[RDM] Relationales Datenbankmodell -- Daten liegen in Tabellen (Relationen) mit Primär- und Fremdschlüsseln, in der dritten Normalform.
  \item[Redundanz] Mehrfaches Speichern derselben Information.
  \item[Referenzielle Integrität] Verhindert verwaiste Datensätze: Ein Fremdschlüssel muss auf einen existierenden Primärschlüssel zeigen.
  \item[Relation] Tabelle im relationalen Modell; eine Zeile ist ein \emph{Tupel}, ein Datensatz ein \emph{Record}.
  \item[Replikation] Dieselben Daten liegen auf mehreren Nodes -- entweder Primary/Secondary oder Peer-to-Peer.
  \item[Sharding] Aufteilen des Datenbestands: Jeder Node hält nur einen Teil der Daten.
  \item[Skalarfunktion] Benutzerdefinierte Funktion, die genau einen Wert zurückgibt (\icode{RETURNS} deklariert den Typ, \icode{RETURN} den Wert).
  \item[Star-Schema] Datenmodell des Data Warehouse: Faktentabelle mit Kennzahlen und Fremdschlüsseln in der Mitte, sternförmig umgeben von bewusst denormalisierten Dimensionstabellen.
  \item[Subselect] Verschachtelte Abfrage; das innere \icode{SELECT} wird zuerst ausgeführt, sein Ergebnis nutzt die äußere Abfrage.
  \item[Teilabhängigkeit] Ein Nichtschlüsselattribut hängt nur von einem Teil eines zusammengesetzten Schlüssels ab -- verletzt die 2NF.
  \item[Transaktion] Logisch zusammenhängende Gruppe von SQL-Befehlen, die vollständig oder gar nicht ausgeführt wird.
  \item[Transitive Abhängigkeit] Ein Nichtschlüsselattribut hängt über ein anderes Nichtschlüsselattribut vom Schlüssel ab ($A \rightarrow C \rightarrow B$) -- verletzt die 3NF.
  \item[Trigger] Vom DBMS automatisch ausgelöster Anweisungsblock -- \icode{BEFORE} oder \icode{AFTER} einem \icode{INSERT}, \icode{UPDATE} oder \icode{DELETE}.
  \item[Tupel] Zeile einer Relation; Liste fester Länge aus Attributwerten.
  \item[UNION] Vereinigt die Ergebnisse zweier \icode{SELECT}-Abfragen zeilenweise (vertikal); entfernt Duplikate, \icode{UNION ALL} behält sie.
  \item[Verteiltes System] Sammlung mehrerer unabhängiger Computer bzw. Prozesse, die nach außen wie ein einziges System wirken.
  \item[View] Gespeicherte \icode{SELECT}-Abfrage, die wie eine virtuelle Tabelle genutzt wird -- ohne eigene Daten.
\end{description}
